\clearpage
\section{Device tree 與 probe：kernel 怎麼找到 driver}

\begin{oneline}
\todo{一句話講完這一節}
\end{oneline}

\begin{before}
\todo{v1.0.0 怎麼做，引用 sw/app/pl\_hal\_fuse.c 的實際程式碼}
\end{before}

\begin{tech}
\todo{用到的 kernel API 與機制}
\end{tech}

\begin{gain}
\todo{具體改善了什麼，能驗證的寫法}
\end{gain}

\subsection{實際例子：程式碼與板上輸出}
\todo{\textbackslash snip 抽出的函式，逐段說明；板上實際 log}

\subsection{原理}
\todo{為什麼這樣做是對的；必要時一張圖}

\subsection{面試會問}
\begin{qa}{device-tree matching：compatible 字串怎麼讓 kernel 找到你的 driver，進而呼叫 probe？}
\todo{30 秒版本的回答}
\end{qa}

\subsection{關鍵函式速查}
\todo{函式、在哪個檔案第幾行、做什麼}
